% !TeX root = ../kaehler-einstein-degeneration.tex
% Current content: 2026-11-20, v1.0.
\section{三重点の小さい積窓と拡大可能な領域}
\label{coarse:sec:triple-windows}

Source を含まない三重点で unit を吸収した exact coordinates を選び，
\begin{equation}\label{coarse:eq:exact-triple-window-model}
 \zeta_0\zeta_1\zeta_2=t,
 \qquad \ell_a=-\log|\zeta_a|^2,
 \qquad L=-\log|t|^2,
 \qquad \sum_{a=0}^2\ell_a=L
\end{equation}
とする．十分大きな固定 $m_0>1$ を取り，
$\min_a\ell_a>m_0$ の部分が一つの enlarged labelled triple chart に含まれるようにする．
以下の領域は常に二つの独立な位相の全軌道を含む．

$\eta_a=d\log\zeta_a$ と置くと $\eta_0+\eta_1+\eta_2=0$ である．
任意の異なる二添字 $j,k$ に対して
\[
 dV_{\log}
  =(\sqrt{-1}\eta_j\wedge\bar\eta_j)
    \wedge(\sqrt{-1}\eta_k\wedge\bar\eta_k)
\]
は同じ正の volume form を与える．比較に用いる volume form と K\"ahler form を
\begin{equation}\label{coarse:eq:triple-cusp-reference}
 \mathcal V_t
  =\frac{8L^2}{\ell_0^2\ell_1^2\ell_2^2}\,dV_{\log},
 \qquad
 \omega_t^{\mathrm{cusp}}
  =\ddc\log\mathcal V_t
  =2\sqrt{-1}\sum_{a=0}^2
       \frac{\eta_a\wedge\bar\eta_a}{\ell_a^2}
\end{equation}
と定める．Volume form の対数への $\ddc$ は局所 holomorphic frame を用いて解釈する．
対応する実計量は
\begin{equation}\label{coarse:eq:constrained-triple-cusp-metric}
 g_t^{\mathrm{cusp}}
 =\left.
    \sum_{a=0}^2\frac{d\ell_a^2+4d\theta_a^2}{\ell_a^2}
   \right|_{\sum\ell_a=L,\ \sum\theta_a=\arg t}.
\end{equation}
この計量には従属座標の rank-one 項が含まれ，一般には二因子の積ではない．
行列式を計算すると
\begin{equation}\label{coarse:eq:triple-cusp-exact-volume-ratio}
 \frac{(\omega_t^{\mathrm{cusp}})^2}{\mathcal V_t}
  =\frac{\ell_0^2+\ell_1^2+\ell_2^2}{L^2}
  \in\left[\frac13,1\right].
\end{equation}

\subsection{Balanced な点での積領域の障害}

\begin{proposition}[一つの積領域には入らない固定半径球]
\label{coarse:prop:triple-product-obstruction}
位相を固定した balanced point
$p_t$ を $\ell_0(p_t)=\ell_1(p_t)=\ell_2(p_t)=L/3$ で定める．
十分大きな $L$ に対し，
$B_{g_t^{\mathrm{cusp}}}(p_t,2)$ を含む二つの独立な algebraic coordinates の
product domain であって，全体が同じ triple chart に含まれるものは存在しない．
ここで product domain は，ある異なる $i,j,k$ に対する
$(\zeta_j,\zeta_k)$ の積領域と，従属座標 $\zeta_i=t/(\zeta_j\zeta_k)$ による
その像を意味する．同じ主張は labelled angular cover 上でも成立する．
\end{proposition}

\begin{proof}
固定した位相の下で，深度が
\[
 (3L/5,L/5,L/5)
\]
の三つの置換となる点を取る．$p_t$ から各点への深度空間内の直線路では
全ての $\ell_a$ が $L/5$ 以上である．その長さは
\[
 5\sqrt{(4/15)^2+2(2/15)^2}
 =\sqrt{8/3}<2
\]
以下であり，経路全体は $L$ が大きければ $\min\ell_a>m_0$ に含まれる．
従って三点は全て所要球に入る．

どの添字 $i$ を従属に選んでも，残る二座標の product domain は，
$\ell_j=3L/5$ の点の $\zeta_j$ 座標と，
$\ell_k=3L/5$ の点の $\zeta_k$ 座標を独立に組み合わせた点を含む．
そこで従属深度は
\[
 \ell_i=L-3L/5-3L/5=-L/5
\]
となる．これは $|\zeta_i|<1$ にすら属さず，triple chart に含まれるという仮定に矛盾する．
角度を持ち上げてもこの絶対値の矛盾は変わらない．
\end{proof}

\begin{remark}\label{coarse:rem:triple-product-obstruction-scope}
この障害は，小さい積窓や各積窓自身の complete metric の存在を否定しない．
否定されるのは，constrained metric の任意の固定半径球を，
二つの独立な大きい annulus windows の積へ一度に含める主張である．
有限アトラスで従属 label を選び直しても，balanced point では同じ三点による障害が残る．
\end{remark}

\subsection{全三重点領域で使える小さい相対窓}

\begin{lemma}[最大深度を従属にする小さい相対積窓]
\label{coarse:lem:triple-relative-window}
$p\in X_t$ が $\min_a\ell_a(p)\ge2m_0$ を満たすとし，
$\ell_a^*=\ell_a(p)$ と書く．$\ell_i^*=\max_a\ell_a^*$ となる添字を選び，
残る二添字を $j,k$ とする．$\delta=1/8$ とし，
\begin{align}
 W_{p,t}
  &=\left\{(\zeta_j,\zeta_k):
       (1-\delta)\ell_b^*<\ell_b<(1+\delta)\ell_b^*
       \ (b=j,k)\right\},
       \label{coarse:eq:triple-relative-window}\\
 W_{p,t}^{\mathrm{in}}
  &=\left\{(\zeta_j,\zeta_k):
       (1-\delta/2)\ell_b^*<\ell_b<(1+\delta/2)\ell_b^*
       \ (b=j,k)\right\}
       \label{coarse:eq:triple-relative-inner-window}
\end{align}
と置く．各領域は全位相を含む二つの有限 annuli の積である．このとき次が成立する．
\begin{enumerate}[label=\textup{(S\arabic*)}]
\item $\zeta_i=t/(\zeta_j\zeta_k)$ により，$\overline{W_{p,t}}$ は同じ enlarged triple chart
内へ正則に写り，その全体で
\begin{equation}\label{coarse:eq:relative-dependent-depth}
 (1-2\delta)\ell_i^*\le\ell_i\le(1+2\delta)\ell_i^*
\end{equation}
が成り立つ．特に $\min_a\ell_a>m_0$ である．
\item $W_{p,t}$ の二つの有限 annuli 自身の complete curvature $-1$ metrics の積を
$\widehat\omega_{p,t}$ とする．$\widehat\omega_{p,t}$ は
$\Ric(\widehat\omega_{p,t})=-\widehat\omega_{p,t}$ を満たし，
ある $t,p$ に依存しない $C>1$ に対し
\begin{equation}\label{coarse:eq:relative-window-tensor-comparison}
 C^{-1}\omega_t^{\mathrm{cusp}}
 \le\widehat\omega_{p,t}\le C\omega_t^{\mathrm{cusp}}
 \qquad\text{on }W_{p,t}^{\mathrm{in}}
\end{equation}
となる．この比較は actual metric に関する仮定を用いない．
\item Actual KE metric は
\begin{equation}\label{coarse:eq:relative-window-volume-upper}
 \omega_t^2\le(\widehat\omega_{p,t})^2
 \quad\text{on }W_{p,t},
 \qquad
 \omega_t^2(p)\le C(\omega_t^{\mathrm{cusp}})^2(p)
\end{equation}
を満たす．
\end{enumerate}
\end{lemma}

\begin{proof}
独立な二深度の変化の総和は
$\delta(\ell_j^*+\ell_k^*)\le2\delta\ell_i^*$ 以下なので，
$\ell_i=L-\ell_j-\ell_k$ から
\eqref{coarse:eq:relative-dependent-depth} を得る．
従属深度は $3\ell_i^*/4$ 以上，独立深度は $7\ell_b^*/8$ 以上であり，
全位相にわたって同じ chart 内に留まる．固定 $t$ では radial closure も
$(\CC^*)^2$ の compact subset である．

$(\eta_j,\eta_k)$ coframe において
$\omega_t^{\mathrm{cusp}}/(2\sqrt{-1})$ の係数行列は
\[
 \begin{pmatrix}\ell_j^{-2}&0\\0&\ell_k^{-2}\end{pmatrix}
 +\ell_i^{-2}
     \begin{pmatrix}1&1\\1&1\end{pmatrix}.
\]
独立深度と中心深度は一様比較可能で，
$\ell_i\ge3\ell_i^*/4$ および $\ell_j^*,\ell_k^*\le\ell_i^*$ である．
従って，$2\sqrt{-1}$ を含めた対応する form は
\[
 2\sqrt{-1}\left(
   \frac{\eta_j\wedge\bar\eta_j}{(\ell_j^*)^2}
   +\frac{\eta_k\wedge\bar\eta_k}{(\ell_k^*)^2}\right)
\]
と $W_{p,t}$ 全体で一様比較可能である．

一方，$A_b=(1-\delta)\ell_b^*$，$B_b=(1+\delta)\ell_b^*$ と置けば，
第 $b$ 因子の intrinsic complete form は
\[
 2\sqrt{-1}
 \left(\frac{\pi}{B_b-A_b}\right)^2
 \frac{\eta_b\wedge\bar\eta_b}
 {\sin^2\!\left(\pi\frac{\ell_b-A_b}{B_b-A_b}\right)}.
\]
内側半窓では sine の引数が $[\pi/4,3\pi/4]$ に入るので，その係数は
$\pi^2/(4\delta^2(\ell_b^*)^2)$ と
$\pi^2/(2\delta^2(\ell_b^*)^2)$ の間にある．
これで \eqref{coarse:eq:relative-window-tensor-comparison} を得る．

最後に \cref{coarse:lem:complete-volume-upper} を適用する．
ここでは角度の compact quotient を用いると，radial closure は固定 $t$ で compact であり，
actual density はその上で有限である．従って
$\log(\omega_t^2/(\widehat\omega_{p,t})^2)$ は全 radial boundary で $-\infty$ へ行き，
最大値を内点で取る．その点で
$\ddc\log(\omega_t^2/(\widehat\omega_{p,t})^2)
=\omega_t-\widehat\omega_{p,t}\le0$ となるから，最大値は $0$ 以下である．
中心 $p\in W_{p,t}^{\mathrm{in}}$ で (S2) を用いれば残る上界を得る．
\end{proof}

\begin{corollary}[全深部三重点領域での volume comparison と有界 potential]
\label{coarse:cor:triple-volume-comparison}
固定した $m_0$ を必要なら大きくすれば，ある $C>1$ に対し
\begin{equation}\label{coarse:eq:full-triple-volume-comparison}
 C^{-1}\mathcal V_t\le\omega_t^2\le C\mathcal V_t
 \qquad\text{on }\{\min_a\ell_a\ge2m_0\}
\end{equation}
が全位相について一様に成り立つ．この領域上で
\begin{equation}\label{coarse:eq:full-triple-bounded-potential}
 u_t=\log\frac{\omega_t^2}{\mathcal V_t},
 \qquad |u_t|\le C,
 \qquad
 \ddc u_t=\omega_t-\omega_t^{\mathrm{cusp}}.
\end{equation}
これらは $\mathrm{(TC)}_5$ を仮定しない．
\end{corollary}

\begin{proof}
上界は各点で \cref{coarse:lem:triple-relative-window} を適用し，
\eqref{coarse:eq:triple-cusp-exact-volume-ratio} を用いれば得られる．
下界には \cref{coarse:cor:global-volume-lower} で証明した
$\omega_t^2\ge c\Omega_t^{\mathrm{low}}$ を用いる．
三つの active labels では exact coordinates と元の defining sections の比が bounded unit であり，
\[
 \alpha_{a,t}=A_0+\min\{\ell_a,L/2\}+O(1)
 \asymp\ell_a.
\]
最後の比較では $2m_0\le\ell_a\le L$ を用いた．
Inactive labels は $O(1)$ であり，canonical density の logarithmic frame も
$dV_{\log}$ と一様比較可能である．従って
$\Omega_t^{\mathrm{low}}\asymp\mathcal V_t$ である．
これで両側 volume bounds を得る．最後の恒等式は
$\ddc\log\omega_t^2=\omega_t$ と
\eqref{coarse:eq:triple-cusp-reference} の差である．
\end{proof}

\begin{remark}\label{coarse:rem:triple-volume-not-tensor}
\eqref{coarse:eq:full-triple-bounded-potential} は本稿の actual compact degeneration に対して
成立する volume と potential の評価であり，tensor comparison とは区別する．
この評価だけから $C^{-1}\omega_t^{\mathrm{cusp}}\le\omega_t\le C\omega_t^{\mathrm{cusp}}$，
actual metric の固定半径球の包含，または actual boundary escape は結論しない．
Bounded density から tensor comparison を導く一般原理への障害は
\cref{coarse:prop:density-counterexample} に示した．
\end{remark}

\subsection{一つの深度が支配する領域での拡大積窓}

\begin{lemma}[Dominant moving sector の product-domain exhaustion]
\label{coarse:lem:moving-product-domain}
$p_t$ の深度を
\[
 (\ell_i^*,\ell_j^*,\ell_k^*)=(L-a_t-b_t,a_t,b_t)
\]
と書き，
\begin{equation}\label{coarse:eq:moving-product-sector-assumptions}
 m_t:=\min\{a_t,b_t\}\longrightarrow\infty,
 \qquad \varepsilon_t:=\frac{a_t+b_t}{L}\longrightarrow0
\end{equation}
を仮定する．$\rho_t\to\infty$ を
\begin{equation}\label{coarse:eq:moving-product-radii}
 e^{\rho_t}\varepsilon_t\longrightarrow0,
 \qquad e^{-\rho_t}m_t\longrightarrow\infty
\end{equation}
となるように選ぶ．例えば
\[
 \rho_t=\frac12\min\left\{
       \log\frac{L}{a_t+b_t},\ \log m_t\right\}
\]
でよい．全位相を含む積領域を
\begin{equation}\label{coarse:eq:moving-product-window}
 W_t^{\mathrm{mov}}
 =\left\{e^{-\rho_t}a_t<\ell_j<e^{\rho_t}a_t,\quad
          e^{-\rho_t}b_t<\ell_k<e^{\rho_t}b_t\right\}
\end{equation}
とする．二因子の cusp product を
\begin{equation}\label{coarse:eq:moving-cusp-product}
 \omega_t^{\mathrm{prod}}
 =2\sqrt{-1}\left(
       \frac{\eta_j\wedge\bar\eta_j}{\ell_j^2}
       +\frac{\eta_k\wedge\bar\eta_k}{\ell_k^2}\right)
\end{equation}
と書き，その実計量を $g_t^{\mathrm{prod}}$ とする．このとき次が成立する．
\begin{enumerate}[label=\textup{(M\arabic*)}]
\item 十分小さい $t$ で $\overline{W_t^{\mathrm{mov}}}$ は同じ triple chart に含まれ，
その全体で $\min_a\ell_a\to\infty$，$\ell_i/L\to1$ となる．
任意の固定 $r>0$ に対し，labelled cover 上の
$B_{g_t^{\mathrm{prod}}}(p_t,r)$ は十分小さい $t$ で $W_t^{\mathrm{mov}}$ に含まれる．
\item $W_t^{\mathrm{mov}}$ 上で
\begin{equation}\label{coarse:eq:moving-cusp-tensor-error}
 0\le\omega_t^{\mathrm{cusp}}-\omega_t^{\mathrm{prod}}
 \le
 \left(\frac{e^{\rho_t}(a_t+b_t)}
 {L-e^{\rho_t}(a_t+b_t)}\right)^2\omega_t^{\mathrm{prod}}.
\end{equation}
また $\mathcal V_t/(\omega_t^{\mathrm{prod}})^2\to1$ も同じ領域で一様に成り立つ．
\item 各 finite annulus 自身の complete curvature $-1$ metric の積を
$\widehat\omega_t^{\mathrm{mov}}$ とする．任意の固定 $r>0$ に対し
\begin{equation}\label{coarse:eq:moving-intrinsic-cusp-comparison}
 \omega_t^{\mathrm{prod}}
 \le\widehat\omega_t^{\mathrm{mov}}
 \le\coth^2\!\left(\frac{\rho_t-r}{2}\right)
                  \omega_t^{\mathrm{prod}}
 \quad\text{on }B_{g_t^{\mathrm{prod}}}(p_t,r)
\end{equation}
が成立する．Normalized holomorphic covering coordinates 上では，
全固定階 $C^k$ で $\widehat\omega_t^{\mathrm{mov}}-\omega_t^{\mathrm{prod}}\to0$ となる．
\item Actual volume は，同じ固定球上で
\begin{equation}\label{coarse:eq:moving-sharp-volume-upper}
 \omega_t^2
 \le\left(\widehat\omega_t^{\mathrm{mov}}\right)^2
 \le\coth^4\!\left(\frac{\rho_t-r}{2}\right)
                  (\omega_t^{\mathrm{prod}})^2
 =(1+o(1))(\omega_t^{\mathrm{prod}})^2
\end{equation}
を満たす．この一側の漸近的上界も $\mathrm{(TC)}_5$ を用いない．
\end{enumerate}
\end{lemma}

\begin{proof}
示した $\rho_t$ は無限大へ行き，
\[
 e^{\rho_t}\varepsilon_t\le\sqrt{\varepsilon_t},
 \qquad
 e^{-\rho_t}m_t\ge\sqrt{m_t}
\]
を満たす．一般に \eqref{coarse:eq:moving-product-radii} の下で
\[
 \ell_j+\ell_k\le e^{\rho_t}(a_t+b_t)=o(L),
 \qquad
 \min\{\ell_j,\ell_k\}\ge e^{-\rho_t}m_t\to\infty.
\]
従って $\ell_i=L-\ell_j-\ell_k=(1-o(1))L$ であり，同じ chart に留まる．
各 cusp factor では $\log\ell_b$ の gradient norm が $1$ なので，
product distance が $r$ 以下なら
$|\log(\ell_j/a_t)|,|\log(\ell_k/b_t)|\le r$ である．
これで固定 product ball の包含を得る．

二つの forms の差は
$2\sqrt{-1}\ell_i^{-2}(\eta_j+\eta_k)\wedge
 (\bar\eta_j+\bar\eta_k)$ である．Cauchy--Schwarz により
\[
 |v_j+v_k|^2
 \le(\ell_j^2+\ell_k^2)
       \left(\frac{|v_j|^2}{\ell_j^2}
             +\frac{|v_k|^2}{\ell_k^2}\right)
\]
なので \eqref{coarse:eq:moving-cusp-tensor-error} が従う．
Volume の正規化は \eqref{coarse:eq:triple-cusp-reference} により
\[
 \frac{\mathcal V_t}{(\omega_t^{\mathrm{prod}})^2}
  =\frac{L^2}{\ell_i^2}\longrightarrow1
\]
と正確に計算できる．

一つの cusp factor の angular universal cover を考える．
その complete metric は $\ell^{-2}(d\ell^2+4d\theta^2)$ である．
$|\log(\ell/a_t)|\le r$ の点から，有限 annulus の二つの radial boundaries
$\ell=e^{\pm\rho_t}a_t$ までの ambient cusp distance は
$\rho_t-r$ 以上である．従ってその点を中心とする半径 $\rho_t-r$ の
ambient hyperbolic ball は finite window に含まれる．
Poincar\'e metric の domain monotonicity と，半径 $s$ の hyperbolic ball
自身の complete metric の中心密度比 $\coth^2(s/2)$ を用いると，
各因子で \eqref{coarse:eq:moving-intrinsic-cusp-comparison} を得る．

高階収束も明示式から従う．第 $j$ 因子で $x=\ell_j/a_t$，
$q_t=e^{-\rho_t}$ と置けば，intrinsic metric と cusp metric の係数比は
\[
 \left(\frac{\pi x}{q_t^{-1}-q_t}\right)^2
 \sin^{-2}\!\left(\pi\frac{x-q_t}{q_t^{-1}-q_t}\right).
\]
任意の固定区間 $e^{-r}\le x\le e^r$ で，この式は $1$ へ全固定階の
$x$-derivatives とともに収束する．Normalized logarithmic holomorphic coordinate
$(\log\zeta_j-\log\zeta_j(p_t))/a_t$ では cusp metric 自身の係数が $t$ に依存しないため，
これが所要 $C^k$ 収束を与える．他方の因子も同様である．

最後に各固定 $t$ で $W_t^{\mathrm{mov}}$ は compact radial closure を持つ annuli の積である．
\Cref{coarse:lem:complete-volume-upper} により
$\omega_t^2\le(\widehat\omega_t^{\mathrm{mov}})^2$ であり，
二因子の metric upper bound の積から
\eqref{coarse:eq:moving-sharp-volume-upper} を得る．
\end{proof}

\begin{corollary}[Open square の一様な保護深度での具体的半径]
\label{coarse:cor:moving-product-eta-radius}
固定 $0<\eta<1/2$ と $0<c<C$ に対し
\[
 cL^\eta\le a_t,b_t\le CL^{1-\eta}
\]
を仮定する．$R_t=\log L$ として
\begin{equation}\label{coarse:eq:moving-product-eta-radius}
 \rho_t=\frac{\eta R_t}{4}
\end{equation}
と選べば，\cref{coarse:lem:moving-product-domain} の全結論が $a_t,b_t$ に一様に成立する．
特に
\[
 e^{\rho_t}\frac{a_t+b_t}{L}\le2CL^{-3\eta/4},
 \qquad
 e^{-\rho_t}\min\{a_t,b_t\}\ge cL^{3\eta/4}.
\]
\end{corollary}

\begin{proof}
$e^{\rho_t}=L^{\eta/4}$ を仮定した二側深度評価へ代入すれば二式を得る．
両式は \eqref{coarse:eq:moving-product-radii} を一様に保証する．
\end{proof}

\begin{remark}\label{coarse:rem:moving-sector-versus-balanced}
この exhaustion は非支配深度が発散し，かつその和が $o(L)$ となる領域に限る．
Balanced region では \cref{coarse:prop:triple-product-obstruction} があるため，
同じ大きい product-domain statement を主張しない．
また (M1) が含めるのは model metric の球であり，actual KE metric の球ではない．
本補題の volume 上界だけから actual tensor comparison は導かない．
Dominant moving sector については，次の \cref{coarse:thm:moving-triple-tensor} で
別途構成する下 barrier と small-potential regularity により actual tensor comparison を証明する．
全閉 corner の比較や，大域的な actual boundary escape までを本補題の結論に含めない．
\end{remark}
